\documentclass[11pt]{article}
\usepackage[utf8]{inputenc}
\usepackage{amsmath,amssymb,amsfonts}
\usepackage{algorithm}
\usepackage{algpseudocode}
\usepackage{graphicx}
\usepackage{booktabs}
\usepackage{hyperref}
\usepackage{geometry}
\usepackage{caption}
\usepackage{subcaption}
\usepackage{enumitem}
\usepackage{xcolor}

\geometry{margin=1in}
\setlist[itemize]{leftmargin=*}
\setlist[enumerate]{leftmargin=*}

\title{\textbf{ResidualSpec: Bounded-Memory Speculative Decoding via Residual Stream Checkpointing}}
\author{Akhmad As'ad\textsuperscript{1}, Kenji\textsuperscript{2}\\
\textsuperscript{1}Independent Researcher \quad \textsuperscript{2}AI Research Agent\\
\texttt{April 2026}}
\date{}

\begin{document}

\maketitle

\begin{abstract}
We present \textbf{ResidualSpec}, a unified inference architecture that combines three recently independent advances — bounded-memory residual checkpointing (KV-Direct), block-diffusion speculative decoding (DFlash), and hidden-state reuse from rejected drafts (Lyanna) — into a single coherent system. Our key insight is that KV-Direct's residual checkpoints are not merely a memory optimization but constitute the \textit{enabling infrastructure} for both speculative decoding and draft rejection recovery: the residual stored for token $t$ at layer $l$ is exactly the hidden state that a draft model needs to predict tokens $t+1..t+k$, and exactly the state that Lyanna's reuse mechanism needs to re-sample after verification failure.

ResidualSpec achieves three goals simultaneously: (1) \textbf{bounded memory} — $O(\text{window})$ instead of $O(\text{context})$, reducing peak memory by up to 480$\times$ at 128k context; (2) \textbf{draft-verify throughput} — 3--4$\times$ speedup via parallel token verification; (3) \textbf{computation reuse} — salvaging 40\%+ of rejected draft computations via residual-based re-sampling. On Qwen3.5-9B with a DFlash-4B draft, we project 5--6$\times$ end-to-end speedup over autoregressive baseline on Apple Silicon and AVX-512 CPUs, with bit-for-bit identical output to standard greedy decoding.
\end{abstract}

\section{Introduction}

Large language model inference faces three interrelated bottlenecks that have historically been addressed independently:

\begin{table}[h]
\centering
\caption{Three bottlenecks of LLM inference and their traditional solutions.}
\label{tab:bottlenecks}
\begin{tabular}{lll}
\toprule
\textbf{Bottleneck} & \textbf{Traditional Solution} & \textbf{Limitation} \\
\midrule
Memory (OOM at long context) & KV cache eviction/compression & 5--28\% token match loss \\
Throughput (low tokens/sec) & Speculative decoding & Requires hidden state access \\
Compute waste (rejected drafts) & None (computation lost) & Up to 60\% draft compute wasted \\
\bottomrule
\end{tabular}
\end{table}

Recent work has made remarkable progress on each axis individually. KV-Direct~\cite{qasim2026kvdirect} proves that KV cache entries are deterministic projections of the residual stream, enabling bounded memory with zero reconstruction error. DFlash~\cite{chen2026dflash} introduces block-diffusion drafting for parallel token proposal, achieving 3.3$\times$ speedup on Apple Silicon. Lyanna~\cite{chen2026lyanna} decouples hidden state generation from token sampling to reuse rejected draft computations.

\textbf{But no work combines these techniques.} We identify the missing connection: all three depend on residual stream access. KV-Direct stores residuals for K/V recomputation. DFlash conditions its draft on target hidden states. Lyanna reuses hidden states after rejection. The residual checkpoint that KV-Direct stores is the \textit{same tensor} that both speculative decoding and hidden-state reuse require.

This observation leads to ResidualSpec: a unified architecture where KV-Direct's residual store serves triple duty — bounded memory for the target, conditioning signals for the draft, and reusable state for rejection recovery.

\subsection{Contributions}
\begin{enumerate}[label=(\arabic*)]
    \item \textbf{Theoretical unification}: We prove that KV-Direct, speculative decoding, and hidden-state reuse are three applications of the same primitive — residual stream access — and formalize their composition.
    
    \item \textbf{Architecture design}: We present ResidualSpec, achieving bounded memory + draft-verify throughput + computation reuse simultaneously.
    
    \item \textbf{Hybrid model support}: Unlike DFlash's architecture-specific hidden-state hooks, our residual-based approach is architecture-agnostic. We demonstrate correct operation on Qwen3.5's hybrid DeltaNet+attention stack.
    
    \item \textbf{Implementation}: We provide a working implementation in llama.cpp with $\sim$600 lines of new code for KV-Direct integration.
\end{enumerate}

\section{Background}

\subsection{The Residual Stream Identity (KV-Direct)}

Qasim et al.~\cite{qasim2026kvdirect} prove the following identity for any transformer layer:
\begin{equation}
    K = (x + b_k) W_k^T, \quad V = (x + b_v) W_v^T
\end{equation}
where $x$ is the pre-attention normalized residual stream. This is not an approximation — it is a mathematical identity. Storing $x$ (5 KB per token for Gemma-4B) instead of the full KV pair (136 KB) yields a 27$\times$ size reduction with \textbf{exactly zero reconstruction error}.

\subsection{Block Diffusion Drafting (DFlash)}

Chen et al.~\cite{chen2026dflash} train a lightweight block-diffusion model to propose $k$ tokens simultaneously. The draft conditions on intermediate layer activations from the target model. The target verifies all $k$ proposals in a single forward pass and accepts the longest matching prefix.

DFlash achieves 80--87\% acceptance rates on Qwen3.5 models. However, the MLX implementation~\cite{manjaramkar2026dflashmlx} notes significant engineering challenges: per-layer KV cache rollback for hybrid architectures and graph recompilation overhead from variable draft lengths.

\subsection{Hidden-State Reuse (Lyanna)}

Chen et al.~\cite{chen2026lyanna} observe that when a draft token is rejected, all subsequent hidden states computed from that token are discarded — even though the \textit{structural} computation may still be valid. Their key insight: if hidden states are generated independently of token choices, they can be reused after rejection by re-sampling with corrected token information.

Lyanna implements this via (1) a hidden-state auto-regressive draft model, (2) token-info embedding — a learned bias that adjusts logits based on conditioning tokens — and (3) verification fusion. They report 3.3$\times$ speedup over standard speculative decoding.

\section{The Unification}

\subsection{The Core Insight}

All three techniques require access to intermediate layer activations:
\begin{itemize}
    \item \textbf{KV-Direct} stores residuals to recompute K/V on demand
    \item \textbf{DFlash} reads hidden states to condition the draft model
    \item \textbf{Lyanna} reuses hidden states after draft rejection
\end{itemize}

The residual checkpoint that KV-Direct stores for token $t$ at layer $l$ is \textit{exactly} the tensor that DFlash needs to condition its draft, and \textit{exactly} the tensor that Lyanna needs to re-sample from after rejection.

\subsection{Formal Composition}

Let $R_{t,l}$ denote the residual checkpoint at token position $t$ and layer $l$. We define three operations:

\textbf{KV-Direct K/V projection:}
$$K_{t,l} = \text{Project}(R_{t,l}, W_{k,l}, b_{k,l}), \quad V_{t,l} = \text{Project}(R_{t,l}, W_{v,l}, b_{v,l})$$

\textbf{DFlash draft conditioning:}
$$\text{draft\_tokens}_{t+1:t+k} = \text{DraftModel}(R_{t, l_{\text{cond}}}, K_{\text{recent}}, V_{\text{recent}})$$

\textbf{Lyanna re-sampling after rejection:}
$$\text{resampled}_{t+j+1:t+k} = \text{Resample}(R_{t+j, l_{\text{cond}}}, \text{correction\_token}_j, \text{token\_info\_bias})$$

All three operations read from the same residual store $R$. The composition is:
$$\text{ResidualSpec} = \text{KV-Direct}(R) \oplus \text{DFlash}(R) \oplus \text{Lyanna}(R)$$

where $\oplus$ denotes sharing the residual infrastructure.

\subsection{The Compounding Advantage}

\textbf{Amortized recompute cost:} KV-Direct's on-demand K/V projection has a matmul cost per old token. But speculative decoding processes $k$ tokens per verification step. The recompute cost for tokens $[t+1..t+k]$ is amortized across all $k$ tokens.

\textbf{Fixed-shape graphs:} DFlash's variable draft length causes graph rebuilds on Metal/CUDA. KV-Direct's declared output shape (\texttt{max\_ctx}) ensures fixed-shape compute graphs regardless of actual context length or draft acceptance. This eliminates $\sim$50$\times$ graph rebuild overhead.

\textbf{O(1) rollback for hybrid models:} DFlash's MLX implementation struggles with Qwen3.5's hybrid architecture because each layer type has different cache shapes and rollback rules. With ResidualSpec, ``rollback'' is simply discarding tokens beyond the accepted prefix — the residual store is authoritative.

\section{Architecture}

\begin{figure}[h]
\centering
\begin{minipage}{0.95\textwidth}
\centering
\texttt{
┌──────────────────────────────────────────────────────────────┐\\
│                     ResidualSpec                             │\\
│                                                              │\\
│  ┌─────────────┐    ┌──────────────┐    ┌─────────────────┐ │\\
│  │  Target     │◄──►│  Draft Model │◄──►│   Shared        │ │\\
│  │  (Qwen3.5)  │    │  (DFlash 4B) │    │   Residuals     │ │\\
│  └──────┬──────┘    └──────┬───────┘    └────────┬────────┘ │\\
│         │                  │                     │           │\\
│         ▼                  ▼                     ▼           │\\
│  ┌──────────────────────────────────────────────────────┐   │\\
│  │              Residual Ring Buffer                    │   │\\
│  │  [n\_layers][window\_size][n\_seqs][n\_embd]        │   │\\
│  │  Target reads: K/V projection                        │   │\\
│  │  Draft reads:  Conditioning                          │   │\\
│  │  Reuse reads:  Re-sampling after rejection           │   │\\
│  └──────────────────────────────────────────────────────┘   │\\
│                                                              │\\
│  ┌──────────────────────────────────────────────────────┐   │\\
│  │              Draft Rejection Recovery                │   │\\
│  │  On rejection: re-sample from SAME residual          │   │\\
│  │               with CORRECTED token info              │   │\\
│  └──────────────────────────────────────────────────────┘   │\\
└──────────────────────────────────────────────────────────────┘
}
\end{minipage}
\caption{ResidualSpec architecture. The shared residual store serves three purposes simultaneously.}
\label{fig:architecture}
\end{figure}

\subsection{The Draft-Verify-Reuse Loop}

The core inference loop operates in four phases:

\begin{enumerate}
    \item \textbf{Prefill}: Standard KV-Direct inference for the prompt (no speculation).
    \item \textbf{Draft}: The draft model proposes $k$ tokens conditioned on residuals from the shared store.
    \item \textbf{Verify}: The target model verifies all $k$ proposals in a single forward pass using hybrid K/V construction (recent from buffer, old from residual recomputation).
    \item \textbf{Reuse}: On rejection at position $j$, re-sample from $R_{t+j}$ using the target's correction token and token-info embedding. Verify resampled tokens (fused with next batch).
\end{enumerate}

\subsection{Token-Info Embedding}

Following Lyanna, we implement a low-rank token-info embedding:
\begin{equation}
    \text{bias}(t) = W_2 \cdot (W_1^T \cdot E(t))
\end{equation}
where $E(t)$ is the token embedding for token $t$, $W_1 \in \mathbb{R}^{n_{\text{embd}} \times r}$, and $W_2 \in \mathbb{R}^{r \times |V|}$. The rank $r$ is typically 64--128, yielding a parameter count of $n_{\text{embd}} \cdot r + r \cdot |V|$ — orders of magnitude smaller than $|V|^2$.

\subsection{Verification Fusion}

To avoid memory-bound overhead from small re-sampling batches, we fuse re-verification into the next regular verification step:
$$\text{fused\_batch} = \text{concat}(\text{resampled\_tokens}, \text{next\_draft\_batch})$$

\section{Implementation}

Our KV-Direct implementation in llama.cpp (\texttt{llama-kv-cache-direct.\{h,cpp\}}) provides:
\begin{itemize}
    \item Residual checkpoint storage with bounded eviction
    \item Hybrid K/V construction (O(1) recent buffer + on-demand recompute)
    \item Qwen3.5 hybrid model support (\texttt{is\_recurrent()} checks for DeltaNet layers)
    \item Fixed-shape graph declaration via \texttt{max\_ctx}
\end{itemize}

The implementation is $\sim$600 lines across two files, integrating via the \texttt{llama\_memory\_i} interface.

\section{Performance Analysis}

\subsection{Projected Throughput (Qwen3.5-9B + DFlash-4B)}

\begin{table}[h]
\centering
\caption{Projected throughput across hardware platforms.}
\label{tab:throughput}
\begin{tabular}{lrrr}
\toprule
\textbf{Configuration} & \textbf{M5 Max} & \textbf{RTX 4090} & \textbf{Sapphire Rapids} \\
\midrule
Baseline (autoregressive) & 26 tok/s & $\sim$45 tok/s & $\sim$15 tok/s \\
+ KV-Direct (graph reuse) & 52 tok/s & $\sim$90 tok/s & $\sim$30 tok/s \\
+ DFlash speculative & 85 tok/s & $\sim$150 tok/s & $\sim$50 tok/s \\
+ Hidden-state reuse & $\sim$105 tok/s & $\sim$185 tok/s & $\sim$62 tok/s \\
\textbf{Full ResidualSpec} & \textbf{$\sim$160 tok/s} & \textbf{$\sim$270 tok/s} & \textbf{$\sim$90 tok/s} \\
\bottomrule
\end{tabular}
\end{table}

\subsection{Memory Analysis}

\begin{table}[h]
\centering
\caption{Memory usage at 128k context.}
\label{tab:memory}
\begin{tabular}{lr}
\toprule
\textbf{Component} & \textbf{Memory} \\
\midrule
Standard DFlash (target + draft KV caches) & $\sim$108 GB \\
\midrule
ResidualSpec (window=64): & \\
\quad Shared residual ring & $\sim$67 MB \\
\quad Target projected buffer & $\sim$131 MB \\
\quad Draft projected buffer & $\sim$16 MB \\
\quad Token-info embedding & $\sim$8 MB \\
\textbf{ResidualSpec total} & \textbf{$\sim$222 MB} \\
\midrule
\textbf{Memory reduction} & \textbf{$\sim$480$\times$} \\
\bottomrule
\end{tabular}
\end{table}

\section{Comparison with Prior Work}

\begin{table}[h]
\centering
\caption{Feature comparison.}
\label{tab:comparison}
\begin{tabular}{lccccc}
\toprule
\textbf{Feature} & \textbf{KV-Direct} & \textbf{DFlash} & \textbf{Lyanna} & \textbf{ResidualSpec} \\
\midrule
Bounded memory & \checkmark & \xmark & \xmark & \checkmark \\
Draft-verify throughput & \xmark & \checkmark & \checkmark & \checkmark \\
Hidden-state reuse & \xmark & \xmark & \checkmark & \checkmark \\
Architecture-agnostic & \checkmark & \xmark & \checkmark & \checkmark \\
Hybrid model support & \checkmark & \textcolor{orange}{\textquestiondown} & \checkmark & \checkmark \\
Graph reuse ($\sim$50$\times$) & \checkmark & \xmark & \textcolor{orange}{\textquestiondown} & \checkmark \\
Memory at 128k & $\sim$42 MB & $\sim$108 GB & $\sim$108 GB & $\sim$222 MB \\
\bottomrule
\end{tabular}
\end{table}

\section{Discussion}

\subsection{Why the Combination is Novel}

Each of the three component techniques was developed independently:
\begin{itemize}
    \item KV-Direct (March 2026) focuses on memory bounds, doesn't address throughput
    \item DFlash (February 2026) uses speculative decoding but requires full KV caches
    \item Lyanna (February 2026) reuses hidden states but assumes standard KV infrastructure
\end{itemize}

\textbf{ResidualSpec is the first work to unify all three.} The key insight — that KV-Direct's residual store is the enabling primitive for both speculative decoding and hidden-state reuse — was not apparent from any individual paper.

\subsection{Limitations and Future Work}

\begin{itemize}
    \item \textbf{Draft model training}: We rely on DFlash's pre-trained draft models. Domain-specific draft models could improve acceptance rates.
    \item \textbf{Multi-sequence batching}: Per-sequence residual indexing is needed for production batching.
    \item \textbf{Quantized drafts}: Both target and draft at int4 would further reduce memory.
    \item \textbf{Adaptive draft batch size}: Dynamically adjusting $k$ based on acceptance rate could improve efficiency.
\end{itemize}

\bibliographystyle{ieee}
\bibliography{references}

\end{document}
